% may27_equations.tex
% Clean, screenshot-friendly typeset of every equation used in the
% "Research Presentation May 27" deck. Each card maps to a slide; the
% intent is to screenshot a card directly onto its slide.
%
% Compile:  pdflatex may27_equations.tex   (run twice for tcolorbox)
%
% Reflects the CURRENT production stack (NOT the May-4 math cards):
%   * counterions K+/SO4^2-  (analytic Bikerman, not ClO4-)
%   * parallel 2e-/4e- ORR:  R2e E0=0.695 V, R4e E0=1.23 V vs RHE
%   * Stern Robin BC, C_S = 0.20 F/m^2
%   * log-c primary unknowns, proton electrochemical potential mu_H
%   * log-rate Butler-Volmer
% Sources: scripts/studies/talk_figs/fig_*.py (the figures on the slides),
%   docs/legacy/PNP Equation Formulations.tex, docs/phase6/singh_2016_pka_formula.md,
%   CLAUDE.md hard rules, writeups/WeekOfMay4/group_meeting_math_cards.tex (style).

\documentclass[11pt]{article}
\usepackage[margin=0.75in]{geometry}
\usepackage{amsmath,amssymb,bm}
\usepackage{xcolor}
\usepackage[most]{tcolorbox}
\usepackage{microtype}
\usepackage{parskip}
\usepackage{siunitx}
\usepackage{enumitem}
\setlist[itemize]{leftmargin=*,topsep=2pt,itemsep=2pt}

% Plain black-on-white card: title on a thin rule above the body.
\tcbset{
  cardstyle/.style={
    enhanced,
    colback=white, colframe=white, coltitle=black,
    fonttitle=\bfseries\large, title={#1},
    boxrule=0pt, frame hidden, arc=0pt,
    left=4pt, right=4pt, top=4pt, bottom=4pt,
    attach boxed title to top left={xshift=0pt, yshift=0pt},
    boxed title style={colback=white, colframe=white, boxrule=0pt, arc=0pt,
      left=0pt, right=0pt, top=0pt, bottom=2pt},
    after title={\par\vspace{2pt}\noindent\rule{\linewidth}{0.4pt}\par\vspace{4pt}},
    breakable=false
  }
}

\newcommand{\VT}{V_{T}}
\newcommand{\Eeq}{E^{\circ}}
\newcommand{\VRHE}{V_{\mathrm{RHE}}}
\newcommand{\CS}{C_{S}}
\newcommand{\dd}{\,\mathrm{d}}
\newcommand{\Jbf}{\mathbf{J}}
\newcommand{\nbf}{\mathbf{n}}

\begin{document}

\begin{center}
  {\Large\bfseries Equation Reference --- Research Presentation (May 27)}\\[3pt]
  {\small Each card maps to a slide; screenshot it directly onto that slide.
  Forms reflect the current production stack: K$^+$/SO$_4^{2-}$ analytic
  Bikerman counterions, parallel 2e$^-$/4e$^-$ ORR, Stern Robin BC
  ($\CS = 0.20$\,F/m$^2$), log-$c$ unknowns with proton electrochemical
  potential $\mu_H$, log-rate Butler--Volmer.}\\[2pt]
  {\footnotesize Symbols: $F$ Faraday constant, $R$ gas constant, $T$
  temperature, $\VT = RT/F$ thermal voltage; $D_i,z_i,c_i$ diffusivity,
  charge, concentration of species $i$; $\varphi$ electrostatic potential;
  $\varepsilon$ permittivity. Dynamic species $\{\mathrm{O_2},
  \mathrm{H_2O_2}, \mathrm{H^+}\}$ ($z=0,0,+1$); counterions
  $\{\mathrm{K^+}, \mathrm{SO_4^{2-}}\}$ carried analytically.}
\end{center}

%% =====================================================================
\begin{tcolorbox}[cardstyle={Slide 5 --- ORR selectivity: two competing pathways}]
We want the 2e$^-$ product (H$_2$O$_2$), not the 4e$^-$ loss to water.
\[
  \text{(2e$^-$, wanted)}\quad
  \mathrm{O_2} + 2\mathrm{H^+} + 2e^- \;\longrightarrow\; \mathrm{H_2O_2},
  \qquad \Eeq = 0.695~\text{V vs RHE}
\]
\[
  \text{(4e$^-$, loss)}\quad
  \mathrm{O_2} + 4\mathrm{H^+} + 4e^- \;\longrightarrow\; 2\,\mathrm{H_2O},
  \qquad \Eeq = 1.23~\text{V vs RHE}
\]
\textbf{Selectivity} = keeping O$_2$ on the 2e$^-$ branch. The two
reactions run \emph{in parallel} at the electrode (Ruggiero 2022 topology).
\end{tcolorbox}

%% =====================================================================
\begin{tcolorbox}[cardstyle={Slide 6 --- RRDE measurement \& transport}]
Rotating ring--disk electrode: the disk makes H$_2$O$_2$, the ring
downstream detects it.
\[
  \text{collection efficiency}\quad N = 0.224 \quad\text{(Ruggiero 2022)}
\]
Rotation sets the diffusion-layer thickness via the Levich relation
(the physical meaning of the domain length $L_{\mathrm{eff}}$):
\[
  L_{\mathrm{eff}} \;\approx\; 1.61\,D_i^{1/3}\,\nu^{1/6}\,\omega^{-1/2},
  \qquad \omega = \frac{2\pi\,\Omega}{60},
\]
with $\nu$ the kinematic viscosity and $\Omega$ the rotation rate (rpm).
For O$_2$ in water this gives $L_{\mathrm{eff}}\approx 86,\,43,\,35~\mu$m
at $\Omega = 400,\,1600,\,2500$ rpm; the model uses $L_{\mathrm{eff}}
\approx 100~\mu$m.
\end{tcolorbox}

%% =====================================================================
\begin{tcolorbox}[cardstyle={Slide 12 --- The PNP--BV model (governing equations + every BC)}]
\textbf{Nernst--Planck} in the diffuse layer, one equation per dynamic
species $i$:
\[
  \frac{\partial c_i}{\partial t}
   = \nabla\!\cdot\!\left(D_i\!\left[\nabla c_i
      + \frac{z_i F}{RT}\,c_i\,\nabla\varphi\right]\right)
   = -\nabla\!\cdot\Jbf_i,
  \qquad
  \Jbf_i = -D_i\!\left(\nabla c_i + \frac{z_i F}{RT}\,c_i\,\nabla\varphi\right).
\]

\textbf{Poisson} for the potential (counterions enter as an analytic
source --- next card):
\[
  -\nabla\!\cdot(\varepsilon\,\nabla\varphi) \;=\; \sum_i z_i F\,c_i .
\]

\textbf{Boundary conditions.}
\begin{itemize}
  \item \emph{Electrode / OHP --- Butler--Volmer flux} (species $i$,
        reactions $r$):
  \[
    -D_i\!\left(\nabla c_i + \frac{z_i F}{RT}\,c_i\,\nabla\varphi\right)\!\cdot\nbf
    \;=\; \sum_r \nu_{i,r}\,R_r,
  \]
  \[
    R_r = k_{0,r}\!\left[\,c_O\,e^{-\alpha_r F\eta_r/RT}
                       \;-\; c_R\,e^{(1-\alpha_r)F\eta_r/RT}\right],
    \qquad
    \eta_r = \varphi_m - \varphi_s - \Eeq_{,r}.
  \]
  Here $c_O,c_R$ are the oxidized / reduced surface concentrations for
  reaction $r$ and $\nu_{i,r}$ the stoichiometric coefficient
  ($\Eeq_{,r}\in\{0.695,\,1.23\}$ V). The log-rate construction of this
  same rate law is on the Slide-14 card.
  \item \emph{Stern compact layer --- Robin BC on $\varphi$:}
  \[
    \varepsilon\,\partial_n\varphi \;=\; \CS\,(V_{\mathrm{app}} - \varphi_{\mathrm{OHP}}),
    \qquad \CS = 0.20~\text{F/m}^2 .
  \]
  \item \emph{Bulk --- Dirichlet:}\quad
  $c_i = c_i^{\mathrm{bulk}}, \quad \varphi = 0.$
\end{itemize}
\end{tcolorbox}

%% =====================================================================
\begin{tcolorbox}[cardstyle={Slide 12 (cont.) --- Analytic Bikerman counterions \& steric crowding}]
Non-reactive ions $b\in\{\mathrm{K^+},\mathrm{SO_4^{2-}}\}$ are \emph{not}
discretized; in local equilibrium with $\varphi$ they enter Poisson
analytically.

\textbf{Ideal (point-ion) Boltzmann limit:}
\[
  c_b(x) \;=\; c_b^{\mathrm{bulk}}\,\exp\!\big(-z_b F\varphi/RT\big).
\]

\textbf{Steric-aware Bikerman closure} (caps packing at $c_b\le 1/a_b$):
\[
  c_b(x) \;=\; c_b^{\mathrm{bulk}}\,
  \frac{\big(1 - \sum_{j\in\mathrm{dyn}} a_j c_j(x)\big)\,
        e^{-z_b F\varphi/RT}}
       {\theta_b + a_b\,c_b^{\mathrm{bulk}}\,e^{-z_b F\varphi/RT}},
  \qquad
  \theta_b = 1 - \!\sum_j a_j c_j^{\mathrm{bulk}} - a_b c_b^{\mathrm{bulk}}.
\]

\textbf{Finite ion size} ($a_i$ = ion volume) adds a steric chemical
potential to every species' driving force:
\[
  \mu_i^{\mathrm{steric}} = -\ln\theta,\quad
  \theta = 1 - \sum_i a_i c_i,\quad
  c_{\max}\approx \tfrac{1}{a}
  \;\;\Rightarrow\;\;
  \Jbf_i = -D_i c_i\big(\nabla u_i + \tfrac{z_iF}{RT}\nabla\varphi
                        + \nabla\mu_i^{\mathrm{steric}}\big).
\]
\end{tcolorbox}

%% =====================================================================
\begin{tcolorbox}[cardstyle={Slide 13 --- Why it is numerically hard}]
In equilibrium each ion piles up (or depletes) Boltzmann-exponentially
at the interface:
\[
  c_i \;\approx\; c_i^{\mathrm{bulk}}\,\exp\!\big(-z_i F\varphi/RT\big)
  \quad\Longrightarrow\quad
  \text{concentrations span} \sim 10\ \text{orders of magnitude.}
\]
The drop occurs across a Debye length that is tiny next to $L_{\mathrm{eff}}$:
\[
  \lambda_D \;=\; \sqrt{\dfrac{\varepsilon RT}{2F^2 I}}\;\ll\; L_{\mathrm{eff}},
  \qquad I = \tfrac12\textstyle\sum_i z_i^2 c_i\ \text{(ionic strength).}
\]
The exponential BV nonlinearity at the boundary feeds this buildup, and
electromigration vs.\ diffusion run on very different timescales ---
together a stiff, ill-conditioned system.
\end{tcolorbox}

%% =====================================================================
\begin{tcolorbox}[cardstyle={Slide 14 --- How we solve it (I): log-$c$ and proton $\mu_H$}]
\textbf{Log-concentration unknowns.} Solve for $u_i=\ln c_i$, not $c_i$;
reconstruct $c_i=e^{u_i}>0$ (positive by construction, tames the buildup):
\[
  u_i = \ln c_i,\qquad
  \Jbf_i = -D_i\,e^{u_i}\!\left(\nabla u_i + \frac{z_i F}{RT}\nabla\varphi\right).
\]

\textbf{Proton electrochemical potential.} For H$^+$ use
\[
  \boxed{\;\mu_H \;=\; u_H + \frac{z_H F}{RT}\,\varphi\;}\quad (z_H=+1),
  \qquad
  c_H = \exp\!\Big(\mu_H - \frac{z_H F}{RT}\varphi\Big),
  \qquad
  \Jbf_H = -D_H\,c_H\,\nabla\mu_H .
\]
In local Boltzmann equilibrium $\nabla\mu_H\approx 0$, so the primary
unknown $\mu_H$ stays \emph{smooth across the double layer} even though
$u_H$ and $\varphi$ each swing by tens of log-units.
\end{tcolorbox}

%% =====================================================================
\begin{tcolorbox}[cardstyle={Slide 14 --- Butler--Volmer: general form vs.\ log-rate construction}]
\textbf{General Butler--Volmer boundary equation} (per reaction $r$; the
Slide-12 form), with $c_O,c_R$ the oxidized / reduced surface
concentrations:
\[
  R_r = k_{0,r}\!\left[\,c_O\,e^{-\alpha_r F\eta_r/RT}
                     \;-\; c_R\,e^{(1-\alpha_r)F\eta_r/RT}\right],
  \qquad \eta_r = \varphi_m - \varphi_s - \Eeq_{,r},
\]
imposed through the flux BC
$-D_i\big(\nabla c_i + \tfrac{z_iF}{RT}c_i\nabla\varphi\big)\!\cdot\nbf
= \sum_r \nu_{i,r} R_r$.

\textbf{Log-rate construction.} With log-concentration unknowns
$c_O = e^{u_O}$, $c_R = e^{u_R}$, fold each concentration \emph{inside}
its exponent, build each branch in log-space, and exponentiate once:
\[
  \log r_r^{(\mathrm{cat})} = \ln k_{0,r} + u_O - \alpha_r\,\tfrac{F\eta_r}{RT},
  \qquad r_r^{(\mathrm{cat})} = \exp\!\big(\log r_r^{(\mathrm{cat})}\big),
\]
\[
  \log r_r^{(\mathrm{anod})} = \ln k_{0,r} + u_R + (1-\alpha_r)\,\tfrac{F\eta_r}{RT},
  \qquad r_r^{(\mathrm{anod})} = \exp\!\big(\log r_r^{(\mathrm{anod})}\big),
\]
\[
  R_r = r_r^{(\mathrm{cat})} - r_r^{(\mathrm{anod})}.
\]

\textbf{The contrast.} Both give an \emph{identical} $R_r$ in exact
arithmetic. The general form evaluates $c_O = \exp(\mathrm{clamp}\,u_O)$
\emph{then} multiplies by the BV exponential, so when $u_O$ is very
negative the clamp pins a finite floor and the rate cannot vanish (a
phantom current). The log-rate form keeps $u_O$ inside the exponent, so
$\log r_r\to-\infty$ smoothly and the branch dies off correctly --- same
equation, numerically robust.
\end{tcolorbox}

%% =====================================================================
\begin{tcolorbox}[cardstyle={Slide 15 --- Continuation: deform easy $\to$ hard along three knobs}]
Never solve the hard problem cold. Ramp three parameters from an easy
regime to the physical target, warm-starting Newton at every step:
\begin{center}
\begin{tabular}{lccl}
\hline
\textbf{knob} & \textbf{start} & \textbf{target} & \textbf{effect}\\
\hline
reaction rate $k_0$        & $10^{-12}$ & $1$ (physical)        & turn the reaction on gradually\\
Stern capacitance $\CS$    & $0.10$     & $0.20$ F/m$^2$        & stiffen the compact layer\\
voltage $\VRHE$            & $-0.5$     & $+1.0$ V              & walk the operating window\\
\hline
\end{tabular}
\end{center}
\[
  \text{solution at step } k+1:\quad
  \mathbf{U}_{k+1} = \text{Newton}\big(\,\mathbf{U}_{k}\ \text{as initial guess}\,\big).
\]
\end{tcolorbox}

%% =====================================================================
\begin{tcolorbox}[cardstyle={Slide 16 --- When a step fails: recover, don't restart}]
If a continuation rung fails between a converged value $a$ and a failed
target $b$, insert the \textbf{geometric ($\sqrt{\ }$) midpoint} and
retry from the last good solution:
\[
  k_{\mathrm{mid}} = \sqrt{a\,b}\,.
\]
Recurse this bisection to walk across the Frumkin cliff (up to depth~5;
8 warm-walk substeps), rebuilding nothing --- only the parameter moves.
\end{tcolorbox}

%% =====================================================================
\begin{tcolorbox}[cardstyle={Slide 17 --- Verification: method of manufactured solutions}]
Pick a smooth exact solution $u^\ast$, insert it into the operator
$\mathcal{L}$ to get a source $f=\mathcal{L}u^\ast$, solve
$\mathcal{L}u_h=f$, and measure the error as the mesh refines ($h\to0$):
\[
  \|u^\ast - u_h\|_{L^2} = \mathcal{O}(h^{p+1}),
  \qquad
  \|u^\ast - u_h\|_{H^1} = \mathcal{O}(h^{p}).
\]
For CG1 elements ($p=1$) the expected slopes are $L^2\approx 2$,
$H^1\approx 1$, and the observed rates match for every field
($u_{\mathrm{O_2}}$, $u_{\mathrm{H_2O_2}}$, $\mu_H$, $\varphi$) on the
production logc\_muh + Cs$^+$/SO$_4^{2-}$ + Stern stack.
\end{tcolorbox}

%% =====================================================================
\begin{tcolorbox}[cardstyle={Slide 18 --- Agreement with prior work (closure reproduction)}]
To compare against the single-reaction closure model we (i) run one
2e$^-$ reaction and (ii) align the open-circuit-potential convention by
shifting the \emph{whole} voltage axis and \emph{both} equilibrium
potentials by the same $\Delta V$:
\[
  \VRHE \to \VRHE - \Delta V,\qquad \Eeq_{,r}\to \Eeq_{,r} - \Delta V,
  \qquad \Delta V = V_{\mathrm{OCP}} = 0.785~\text{V (pH 2)}.
\]
The overpotential is invariant, so kinetics are unchanged:
\[
  \eta_r = (\VRHE-\Delta V) - (\Eeq_{,r}-\Delta V) = \VRHE - \Eeq_{,r},
  \qquad
  V_{\mathrm{OCP}} = 0.47 + 0.197 + 0.059\,\mathrm{pH}.
\]
Profiles are reconstructed from the muh unknown:
$c_H = \exp\!\big(\mu_H - \tfrac{z_H F}{RT}\varphi\big)$.
\textbf{Result:} 25/25 converge; O$_2$-ratio and profiles within $\sim5\%$.
\end{tcolorbox}

%% =====================================================================
\begin{tcolorbox}[cardstyle={Slide 20 --- Tentative: cation hydrolysis (field-shifted p$K_a$)}]
A polarized interface turns the hydrated cation into a \emph{local proton
source} --- reachable by the general solver, not the closure:
\[
  \mathrm{M^+(H_2O)}_n + \mathrm{H_2O}
  \;\rightleftharpoons\;
  \mathrm{MOH^0(H_2O)}_{n-1} + \mathrm{H_3O^+}.
\]
The acidity shifts with the interfacial field (Singh 2016, Eq.~3--4):
\[
  \mathrm{p}K_a(\sigma) = \mathrm{p}K_a^{\mathrm{bulk}} + \Delta \mathrm{p}K_a(\sigma),
  \qquad
  \mathrm{p}K_a^{\mathrm{bulk}} = B - \frac{A\,z^2}{r_{M\text{-}O}},
\]
\[
  \Delta \mathrm{p}K_a(\sigma)
   = 2A\,z\,\sigma\,r_{H\text{-}El}\!\left(1 - \frac{r_{M\text{-}O}^2}{r_{H\text{-}El}^2}\right),
  \qquad
  \sigma = \CS\,(\varphi_m - \varphi_s)\ \text{(local Stern charge)}.
\]
Constants $A=620.32$ pm, $B=17.154$, $r_O=63$ pm; the new surface proton
flux is gated by a Langmuir-capped cation coverage $\Gamma$.
\end{tcolorbox}

\end{document}
